\chapter{Metatheory: strong typing, uniqueness, Church--Rosser}
\label{chap:metatheory}

\section*{At a glance}
\begin{itemize}
\item \lead{Goal} The metatheory of \texttt{VEnv.IsDefEq}: regularity, unique typing,
      confluence, syntax-directed typing.
\item \lead{Main results} Theorem~\ref{thm:unique-typing} (unique typing);
      Theorem~\ref{thm:parreds-church-rosser} (Church--Rosser); Theorem~\ref{thm:parreds-standard}
      (standardization); Theorem~\ref{thm:infertype-exists} (inference completeness).
\item \lead{Status} \stSorry{} Theorem~\ref{thm:isdefequ-weakn-iff},
      Theorem~\ref{thm:normaleq-parred}; \stTainted{} unique typing, Church--Rosser,
      standardization, inference; \stProved{} \texttt{Strong.lean}, taking
      \texttt{OrderedStrong} as a hypothesis rather than deriving it.
\item \lead{Contribution} \texttt{iota}: the strong $\iota$ rule
      (Definition~\ref{inductive:isdefeqstrong-pat}); a \texttt{pat} case in eleven inductions
      of \texttt{Strong.lean}; \texttt{OrderedStrong}
      (Definition~\ref{struct:orderedstrong}); one class field/rule case in
      \texttt{ChurchRosser.lean}.
\item \lead{Lean files} \texttt{Lean4Lean/Theory/Typing/\{Strong,UniqueTyping,ChurchRosser,HeadReduction\}.lean}.
\item \lead{Thesis} typesys.tex, unique.tex; standardization: Kashima (2000).
\end{itemize}

\begin{remark}
\lead{Notation} In \texttt{ChurchRosser.lean} and \texttt{HeadReduction.lean}, the reduction
relations below fix an environment/universe count from class \texttt{Params}
(Definition~\ref{class:cr-params}).
\par\smallskip\noindent
\begin{tabular}{p{0.31\linewidth}p{0.62\linewidth}}
\toprule
Notation & Meaning \\
\midrule
$\Gamma \vdash e_1 \equiv e_2 : A$, $\Gamma \vdash e : A$ & \texttt{IsDefEq}, its diagonal
\texttt{HasType} \\
$\Gamma \vdash e_1 \equiv e_2$ & untyped \texttt{IsDefEqU} \\
$\Gamma \Vdash e_1 \equiv e_2 : A$ & \texttt{IsDefEqStrong}, fully annotated \\
$\Gamma \Vdash e : A$, $\Gamma \Vdash e :!\,A$ & \texttt{HasTypeStrong}, with/without conversion \\
$\Gamma \vdash_n e : A$ & stratified judgment \\
$\equiv_p$, $\gg$, $\ggg$, $\gg^*$, $\gg\ll$ & normal equality; parallel, complete, iterated
parallel reduction; Church--Rosser equality \\
$\rightsquigarrow$, $\rightsquigarrow^*$, $\rightsquigarrow_<$ & weak head reduction, iterated,
standard \\
$\triangleright$, $\triangleright^*$ & type inference, up to head reduction \\
\bottomrule
\end{tabular}\par\smallskip
\end{remark}

\section{The strong judgments}

\begin{definition}[Strong def-eq \texttt{IsDefEqStrong}]
\label{def:isdefeqstrong}
\lean{Lean4Lean.VEnv.IsDefEqStrong}
\leanok
\uses{def:isdefeq, def:lookup, struct:venv, def:pattern-check-realizes}
\contrib{mixed}\srcloc{Lean4Lean/Theory/Typing/Strong.lean}{18}
\thesisref{typesys.tex thm:reg}

\lead{In short} A copy of \texttt{IsDefEq} where every rule carries \hl{the typing derivations
its statement needs}.
\begin{itemize}
\item \texttt{bvar}: a sort derivation for the looked-up type;
\item \texttt{appDF}: types of domain, codomain, instantiated codomain;
\item \texttt{beta}: typing of the reduct $e[e']$; \texttt{eta}: the lifted typings;
\item \texttt{extra}: both sides of the axiom, empty and current context;
\item \texttt{iota} adds one constructor, \texttt{pat}
      (Definition~\ref{inductive:isdefeqstrong-pat}).
\end{itemize}
\end{definition}

\begin{definition}[Strong $\iota$ rule]
\label{inductive:isdefeqstrong-pat}
\lean{Lean4Lean.VEnv.IsDefEqStrong.pat}
\leanok
\uses{rule:isdefeq-pat, def:isdefeqstrong, def:pattern-check-realizes, def:patstrong}
\contrib{iota}\srcloc{Lean4Lean/Theory/Typing/Strong.lean}{89}
\thesisref{typesys.tex item:red\_equiv}

\lead{In short} Strengthens \texttt{IsDefEq.pat} with \hl{one extra premise: the reduct is
itself well typed}.
\begin{itemize}
\item \lead{Given} $env.pats\;p\;r$, $p$ matches $e$ at $m_1, m_2$, $\Gamma \Vdash e : A$,
      $\Gamma \Vdash r_1.apply\;m_1\;m_2 : A$, and a check list $chk$ realizing $r_2$ with
      $\Gamma \Vdash t_1 \equiv t_{2.1} : t_{2.2}$ for every $t \in chk$;
\item \lead{Then} $\Gamma \Vdash e \equiv r_1.apply\;m_1\;m_2 : A$.
\end{itemize}
The recorded reduct typing lets five later \texttt{pat} cases close.
\end{definition}

\begin{definition}[Strong typing \texttt{HasTypeStrong}]
\label{def:hastypestrong}
\lean{Lean4Lean.VEnv.HasTypeStrong}
\leanok
\uses{def:isdefeqstrong}
\srcloc{Lean4Lean/Theory/Typing/Strong.lean}{108}

\lead{In short} $\Gamma \Vdash e :!\,A$ for the structural rules, $\Gamma \Vdash e : A$ after
\texttt{base}/\texttt{defeq}, \hl{indexed by a Boolean so inversion on a structural rule is one
\texttt{cases}}. No $\iota$ rule: it types, it does not equate.
\end{definition}

\begin{lemma}[Basic strong operations]
\label{fam:strong-basic}
\lean{Lean4Lean.VEnv.IsDefEqStrong.hasType,Lean4Lean.VEnv.IsDefEqStrong.weak0,Lean4Lean.VEnv.IsDefEqStrong.defeqDF_l,Lean4Lean.VEnv.IsDefEqStrong.instDF}
\leanok
\srcloc{Lean4Lean/Theory/Typing/Strong.lean}{147}

\lead{In short} Basic closure properties.
\begin{itemize}
\item reflexivity of both sides; weaken a closed judgment into any context;
\item retype along a definitionally equal context head;
\item substitution congruence: $A::\Gamma \Vdash f\equiv f':B$ and $\Gamma\Vdash a\equiv a':A$
      give $\Gamma\Vdash f[a]\equiv f'[a']:B[a]$.
\end{itemize}
\end{lemma}
\begin{proof}
\uses{def:isdefeqstrong, thm:strong-instn, thm:strong-weakn}
\leanok
\lead{Idea} \texttt{hasType} is $H.trans\,H.symm$; \texttt{weak0} lifts from the empty context;
\texttt{instDF} builds and contracts a $\beta$ redex, applying the rule at $B$ and at
$\mathsf{sort}\,v$.
\end{proof}

\begin{lemma}[Weakening, strong judgment]
\label{thm:strong-weakn}
\lean{Lean4Lean.VEnv.IsDefEqStrong.weakN}
\leanok
\contrib{mixed}\srcloc{Lean4Lean/Theory/Typing/Strong.lean}{153}

\lead{In short} For ordered $env$ and $W:\mathtt{Ctx.LiftN}\,n\,k\,\Gamma\,\Gamma'$,
$\Gamma \Vdash e_1\equiv e_2:A$ gives
$\Gamma'\Vdash e_1{\uparrow}\equiv e_2{\uparrow}:A{\uparrow}$.
\end{lemma}
\begin{proof}
\uses{def:isdefeqstrong, family:pattern-transport-master, family:pattern-transport-iota}
\leanok
\lead{Idea} Induction; the \texttt{iota} case rewrites with \texttt{Pattern.RHS.liftN\_apply},
lifts the match and realizer, and lifts each check triple pointwise.
\thesisref{typesys.tex thm:weak}
\end{proof}

\begin{theorem}[Strong implies ordinary]
\label{thm:strong-defeq}
\lean{Lean4Lean.VEnv.IsDefEqStrong.defeq}
\leanok
\contrib{mixed}\srcloc{Lean4Lean/Theory/Typing/Strong.lean}{201}

\lead{In short} $\Gamma\Vdash e_1\equiv e_2:A$ implies $\Gamma\vdash e_1\equiv e_2:A$:
\hl{forgetting the annotations is sound}.
\end{theorem}
\begin{proof}
\uses{def:isdefeq, def:isdefeqstrong, rule:isdefeq-pat}
\leanok
\lead{Idea} Induction, applying the matching \texttt{IsDefEq} constructor each time; the
\texttt{pat} case drops the reduct-typing annotation.
\thesisref{no thesis counterpart}
\end{proof}

\begin{lemma}[Monotonicity in environment]
\label{thm:strong-mono}
\lean{Lean4Lean.VEnv.IsDefEqStrong.mono}
\leanok
\contrib{mixed}\srcloc{Lean4Lean/Theory/Typing/Strong.lean}{219}

\lead{In short} $env\le env'$ carries $\Gamma\Vdash e_1\equiv e_2:A$ over $env$ to over $env'$.
\end{lemma}
\begin{proof}
\uses{def:isdefeqstrong, struct:venv-le}
\leanok
\lead{Idea} Induction; \texttt{constDF}/\texttt{extra} use the first two components of the
ordering, the \texttt{pat} case the third, monotonicity of \texttt{pats}.
\thesisref{no thesis counterpart}
\end{proof}

\begin{definition}[Equality up to levels]
\label{def:equptolevels}
\lean{Lean4Lean.VEnv.EqUpToLevels,Lean4Lean.VEnv.EqUpToLevels.instL,Lean4Lean.VEnv.EqUpToLevels.refl,Lean4Lean.VEnv.EqUpToLevels.weakN,Lean4Lean.VEnv.EqUpToLevels.instN}
\leanok
\uses{def:isdefeqstrong, ind:vlevel}
\contrib{mixed}\srcloc{Lean4Lean/Theory/Typing/Strong.lean}{246}
\thesisref{no thesis counterpart}

\lead{In short} Structural congruence relating terms that differ only in \hl{equivalent
universe levels} ($\approx$ on \texttt{const} level lists and \texttt{sort}).
\begin{itemize}
\item two level instantiations of a strongly typed term are related;
\item such a term relates to itself; the relation survives weakening and substitution.
\end{itemize}
\end{definition}

\begin{definition}[Strong invariant \texttt{EnvStrong}]
\label{def:envstrong}
\lean{Lean4Lean.VEnv.EnvStrong}
\leanok
\uses{def:isdefeqstrong}
\srcloc{Lean4Lean/Theory/Typing/Strong.lean}{309}

\lead{In short} $EnvStrong\;env\;U\;e\;A$: \hl{the four invariants regularity needs} about a
closed, well-typed declaration.
\begin{itemize}
\item $[\,]\Vdash e\equiv e:A$;
\item $A$ is strongly a type in the empty context;
\item equivalent level instantiations of $e$ are strongly equal at $A$, left-instantiated;
\item if $e=\Pi A'.B'$ then $A'$ is strongly a type and $B'$ is strongly a type in $[A']$.
\end{itemize}
\end{definition}

\begin{lemma}[\texttt{EnvStrong} is monotone]
\label{thm:envstrong-mono}
\lean{Lean4Lean.VEnv.EnvStrong.mono}
\leanok
\contrib{iota}\srcloc{Lean4Lean/Theory/Typing/Strong.lean}{319}

\lead{In short} $env\le env'$ carries $EnvStrong\;env\;U\;e\;A$ to $env'$.
\end{lemma}
\begin{proof}
\uses{thm:strong-mono, def:envstrong}
\leanok
\lead{Idea} Componentwise, by \texttt{IsDefEqStrong.mono} on each of the four conjuncts.
\thesisref{no thesis counterpart}
\end{proof}

\begin{lemma}[Level instantiation, strong]
\label{thm:strong-instl}
\lean{Lean4Lean.VEnv.IsDefEqStrong.instL}
\leanok
\contrib{mixed}\srcloc{Lean4Lean/Theory/Typing/Strong.lean}{326}

\lead{In short} For levels $ls$ well formed in $U'$: $\Gamma\Vdash e_1\equiv e_2:A$ over $U$
gives $\Gamma[ls]\Vdash e_1[ls]\equiv e_2[ls]:A[ls]$ over $U'$.
\end{lemma}
\begin{proof}
\uses{def:isdefeqstrong, family:pattern-transport-iota}
\leanok
\lead{Idea} Induction; the \texttt{pat} case rewrites with \texttt{Pattern.RHS.instL\_apply},
transports the match and realizer, instantiates each check triple.
\thesisref{no thesis counterpart}
\end{proof}

\begin{definition}[Strong context \texttt{CtxStrong}]
\label{def:ctxstrong}
\lean{Lean4Lean.VEnv.CtxStrong,Lean4Lean.VEnv.CtxStrong.lookup,Lean4Lean.VEnv.CtxStrong.defeq,Lean4Lean.VEnv.CtxStrong.levelWF}
\leanok
\uses{def:isdefeqstrong, def:onctx}
\srcloc{Lean4Lean/Theory/Typing/Strong.lean}{369}

\lead{In short} $CtxStrong\;env\;U\;\Gamma$: \hl{every entry of $\Gamma$ is strongly a sort in
its own prefix}. Gives the lifted type of a de~Bruijn lookup, ordinary well-formedness
$OnCtx\;\Gamma\;(env.IsType\;U)$, and level well-formedness of the entries.
\end{definition}

\begin{theorem}[Substitution, strong judgment]
\label{thm:strong-instn}
\lean{Lean4Lean.VEnv.IsDefEqStrong.instN}
\leanok
\contrib{mixed}\srcloc{Lean4Lean/Theory/Typing/Strong.lean}{385}

\lead{In short} For a strong $\Gamma_0\Vdash e_0\equiv e_0:A_0$, a cut
$W:\mathtt{Ctx.InstN}\,\Gamma_0\,e_0\,A_0\,k\,\Gamma_1\,\Gamma$ and strong $\Gamma$:
$\Gamma_1\Vdash e_1\equiv e_2:A$ gives
$\Gamma\Vdash e_1[e_0]_k\equiv e_2[e_0]_k:A[e_0]_k$.
\end{theorem}
\begin{proof}
\uses{def:isdefeqstrong, family:pattern-transport-master, family:pattern-transport-iota}
\leanok
\lead{Idea} Induction on the derivation, inner induction on $W$ at \texttt{bvar}; the
\texttt{iota} case mirrors weakening with \texttt{Pattern.RHS.instN\_apply},
\texttt{Pattern.matches\_instN}, \texttt{Realizes.map\_instN}.
\thesisref{typesys.tex thm:subst}
\end{proof}

\section{Regularity, read off the annotations}

\begin{theorem}[$\Pi$-inversion, strong]
\label{thm:strong-foralle-inv}
\lean{Lean4Lean.VEnv.IsDefEqStrong.forallE_inv'}
\leanok
\contrib{mixed}\srcloc{Lean4Lean/Theory/Typing/Strong.lean}{472}

\lead{In short} Over an ordered environment strong on constants/axioms and a strong context: if
$\Gamma\Vdash e_1\equiv e_2:V$ and \hl{either side is $\Pi A.B$}, then $A$ is strongly a type in
$\Gamma$ and $B$ is strongly a type in $A::\Gamma$.
\end{theorem}
\begin{proof}
\uses{def:isdefeqstrong, def:envstrong, thm:strong-instn}
\leanok
\lead{Idea} Induction, carrying the equation as a disjunction over the two sides; the
\texttt{iota} case recurses into whichever annotation (redex or reduct) matches the $\Pi$ side.
\thesisref{typesys.tex thm:reg}
\end{proof}

\begin{theorem}[Regularity, strong judgment]
\label{thm:strong-istype}
\lean{Lean4Lean.VEnv.IsDefEqStrong.isType'}
\leanok
\contrib{mixed}\srcloc{Lean4Lean/Theory/Typing/Strong.lean}{525}

\lead{In short} Under the same hypotheses, $\Gamma\Vdash e_1\equiv e_2:A$ implies $A$ is
strongly a type: $\exists u,\ \Gamma\Vdash A\equiv A:\mathsf{sort}\,u$.
\end{theorem}
\begin{proof}
\uses{def:isdefeqstrong, def:envstrong, def:ctxstrong, thm:strong-foralle-inv}
\leanok
\lead{Idea} Induction, reading the annotation off each rule; \texttt{appDF}/\texttt{beta}
substitute the codomain with \texttt{instN}, \texttt{constDF}/\texttt{extra} use the
environment invariant, and the \texttt{pat} case returns the redex's typing annotation.
\thesisref{typesys.tex thm:reg}
\end{proof}

\begin{lemma}[Level-equal replacement]
\label{thm:equptolevels-defeq}
\lean{Lean4Lean.VEnv.EqUpToLevels.defeq}
\leanok
\contrib{mixed}\srcloc{Lean4Lean/Theory/Typing/Strong.lean}{577}

\lead{In short} Under the same environment/context hypotheses: if $\Gamma\Vdash e_1\equiv e_2:A$
and $e_1$ is level-equal to $e_1'$, $e_2$ to $e_2'$, then $\Gamma\Vdash e_1'\equiv e_2':A$.
\hl{A transport lemma, not a claim that level-equal terms are equal on their own.}
\end{lemma}
\begin{proof}
\uses{def:equptolevels, def:isdefeqstrong}
\leanok
\lead{Idea} Induction with inversion on the two level-equalities; the \texttt{iota} case
sandwiches the unchanged rule instance between the two induction hypotheses by transitivity.
\thesisref{no thesis counterpart}
\end{proof}

\section{The environment hypothesis and strengthening}

\begin{definition}[\texttt{PatStrong}]
\label{def:patstrong}
\lean{Lean4Lean.VEnv.PatStrong}
\leanok
\uses{def:isdefeqstrong, def:ctxstrong, def:pattern-check-realizes}
\contrib{iota}\srcloc{Lean4Lean/Theory/Typing/Strong.lean}{663}
\thesisref{typesys.tex item:red\_equiv; unique.tex (1)}

\lead{In short} $PatStrong\;env\;p\;r$: \hl{subject reduction for one rule}, so the obligation
can be discharged rule by rule.
\begin{itemize}
\item \lead{Given} every $U$, strong $\Gamma$, $e$ matching $p$ at $m_1,m_2$ with
      $\Gamma\Vdash e\equiv e:A$, a realizer $chk$ of $r_2$, each triple of $chk$ a strong
      equality;
\item \lead{Then} $\Gamma\Vdash r_1.apply\;m_1\;m_2\equiv r_1.apply\;m_1\;m_2:A$.
\end{itemize}
Used once, at \srcloc{Lean4Lean/Theory/Typing/Strong.lean}{747}.
\end{definition}

\begin{definition}[\texttt{PatsStrongOn}]
\label{def:patsstrongon}
\lean{Lean4Lean.VEnv.PatsStrongOn}
\leanok
\uses{def:patstrong, struct:venv, def:venv-addpat}
\contrib{iota}\srcloc{Lean4Lean/Theory/Typing/Strong.lean}{672}
\thesisref{unique.tex (1)}

\lead{In short} $PatsStrongOn\;env$: every registered rule $(p,r)\in env.pats$ satisfies
$PatStrong\;env\;p\;r$.
\lead{Why} Not a consequence of \texttt{Ordered}, which allows arbitrary well-typed
definitional axioms: under one, an $\iota$ redex can \hl{be well typed while its reduct is
not} (concrete counterexample: \srcloc{Lean4Lean/Theory/Typing/EnvLemmas.lean}{327}).
\end{definition}

\begin{definition}[\texttt{OrderedStrong}]
\label{struct:orderedstrong}
\lean{Lean4Lean.VEnv.OrderedStrong,Lean4Lean.VEnv.OrderedStrong.ordered,Lean4Lean.VEnv.OrderedStrong.strong,Lean4Lean.VEnv.OrderedStrong.pats}
\leanok
\uses{def:ordered, def:envstrong, def:patsstrongon, def:ontypes}
\contrib{iota}\srcloc{Lean4Lean/Theory/Typing/Strong.lean}{679}
\thesisref{no thesis counterpart}

\lead{In short} Bundles $Ordered\;env$, $OnTypes\;env\;(EnvStrong\;env)$ (every constant type
and both sides of every axiom strongly typed) and $PatsStrongOn\;env$ -- \hl{the hypothesis
under which any \texttt{IsDefEq} derivation strengthens}.
\lead{Caveat} Reachable from \texttt{VEnv.WF} only through a \stSorry{} obligation
(Theorem~\ref{thm:wf-patsstrong}, \S Caveats): the chapter's main liability.
\end{definition}

\begin{definition}[\texttt{OrderedStrong} to \texttt{Ordered}]
\label{instance:coeout-orderedstrong}
\lean{Lean4Lean.VEnv.instCoeOutOrderedStrongOrdered}
\leanok
\uses{struct:orderedstrong, def:ordered}
\contrib{iota}\srcloc{Lean4Lean/Theory/Typing/Strong.lean}{684}
\thesisref{no thesis counterpart}

\lead{In short} An anonymous \texttt{CoeOut} instance letting an \texttt{OrderedStrong}
hypothesis stand where \texttt{Ordered} is expected.
\lead{Caveat} With its sibling \texttt{CoeOut (VEnv.WF env) env.OrderedStrong}
(\srcloc{Lean4Lean/Theory/Typing/EnvLemmas.lean}{343}), this hides the \stSorry{}-backed
hypothesis at call sites.
\end{definition}

\begin{theorem}[Strengthening, parameterized]
\label{thm:isdefeq-strong-prime}
\lean{Lean4Lean.VEnv.IsDefEq.strong'}
\leanok
\contrib{mixed}\srcloc{Lean4Lean/Theory/Typing/Strong.lean}{688}

\lead{In short} For $henv:Ordered\;env$, $envIH:OnTypes\;env\;(EnvStrong\;env)$,
$hpats:PatsStrongOn\;env$ and a strong $\Gamma$: every $\Gamma\vdash e_1\equiv e_2:A$ has a
strong derivation $\Gamma\Vdash e_1\equiv e_2:A$.
\end{theorem}
\begin{proof}
\uses{def:isdefeqstrong, def:patsstrongon, thm:strong-istype, thm:strong-foralle-inv, fam:strong-basic}
\leanok
\lead{Idea} Induction, reconstructing each missing annotation; the \texttt{iota} case
strengthens the redex typing and check triples, then applies $hpats$ for the reduct typing.
\thesisref{typesys.tex thm:reg}
\end{proof}

\begin{lemma}[Context strongly well-formed]
\label{thm:ctxstrong-strong-prime}
\lean{Lean4Lean.VEnv.CtxStrong.strong'}
\leanok
\contrib{mixed}\srcloc{Lean4Lean/Theory/Typing/Strong.lean}{749}

\lead{In short} Under the same three environment hypotheses,
$OnCtx\;\Gamma\;(env.IsType\;U)$ implies $CtxStrong\;env\;U\;\Gamma$: \hl{the context half of
Regularity}.
\end{lemma}
\begin{proof}
\uses{thm:isdefeq-strong-prime, def:ctxstrong}
\leanok
\lead{Idea} Induction on the context, strengthening each entry with \texttt{IsDefEq.strong'} in
its own prefix.
\thesisref{no thesis counterpart}
\end{proof}

\begin{theorem}[Closed judgment satisfies \texttt{EnvStrong}]
\label{thm:envstrong-of-hastype}
\lean{Lean4Lean.VEnv.EnvStrong.of_hasType}
\leanok
\contrib{iota}\srcloc{Lean4Lean/Theory/Typing/Strong.lean}{759}

\lead{In short} If $env$ is ordered, all constants/axioms are strongly typed and every
registered rule has subject reduction, then $env\vdash e:A$ in the empty context gives
$EnvStrong\;env\;U\;e\;A$.
\end{theorem}
\begin{proof}
\uses{thm:isdefeq-strong-prime, thm:strong-istype, thm:equptolevels-defeq, thm:strong-foralle-inv, def:envstrong}
\leanok
\lead{Idea} Assemble the four \texttt{EnvStrong} components: \texttt{IsDefEq.strong'},
\texttt{isType'}, \texttt{EqUpToLevels.defeq}+\texttt{instL}, \texttt{forallE\_inv'}.
\thesisref{no thesis counterpart}
\end{proof}

\begin{lemma}[Adding a constant]
\label{thm:ontypes-addconst}
\lean{Lean4Lean.VEnv.OnTypes.addConst}
\leanok
\contrib{iota}\srcloc{Lean4Lean/Theory/Typing/Strong.lean}{770}

\lead{In short} If $env$ is ordered, strong and has subject reduction, $ci$ is well formed and
$env.addConst\;n\;ci=\mathsf{some}\;env'$, then $OnTypes\;env'\;(EnvStrong\;env')$.
\end{lemma}
\begin{proof}
\uses{thm:envstrong-of-hastype, thm:envstrong-mono, def:venv-basic}
\leanok
\lead{Idea} Transport the old invariant along \texttt{addConst\_le} with
\texttt{EnvStrong.mono}, then discharge the new constant with \texttt{EnvStrong.of\_hasType}.
\thesisref{no thesis counterpart}
\end{proof}

\begin{lemma}[Adding an axiom]
\label{thm:ontypes-adddefeq}
\lean{Lean4Lean.VEnv.OnTypes.addDefEq}
\leanok
\contrib{iota}\srcloc{Lean4Lean/Theory/Typing/Strong.lean}{785}

\lead{In short} If $env$ is strong and both sides of $df$ satisfy \texttt{EnvStrong}, so does
$env.addDefEq\;df$.
\end{lemma}
\begin{proof}
\uses{thm:envstrong-mono, def:venv-basic}
\leanok
\lead{Idea} Monotonicity along \texttt{addDefEq\_le}, then case split on whether the axiom is
the new one.
\thesisref{no thesis counterpart}
\end{proof}

\begin{lemma}[Registering a rule]
\label{thm:ontypes-addpat}
\lean{Lean4Lean.VEnv.OnTypes.addPat}
\leanok
\contrib{iota}\srcloc{Lean4Lean/Theory/Typing/Strong.lean}{797}

\lead{In short} $env.addPat\;p\;r$ is strong whenever $env$ is: registering a rule adds neither
a constant nor an axiom.
\end{lemma}
\begin{proof}
\uses{thm:envstrong-mono, fam:addpat, def:venv-addpat}
\leanok
\lead{Idea} Monotonicity of \texttt{EnvStrong} along \texttt{addPat\_le}.
\lead{Caveat} Says nothing about subject reduction of the new rule -- the deferred
\texttt{patsStrong} obligation.
\thesisref{no thesis counterpart}
\end{proof}

\begin{theorem}[Strengthening, strong environment]
\label{thm:isdefeq-strong}
\lean{Lean4Lean.VEnv.IsDefEq.strong,Lean4Lean.VEnv.CtxStrong.strong}
\leanok
\contrib{mixed}\srcloc{Lean4Lean/Theory/Typing/Strong.lean}{805}

\lead{In short} For $OrderedStrong\;env$ and $OnCtx\;\Gamma\;(env.IsType\;U)$: \hl{every
$\Gamma\vdash e_1\equiv e_2:A$ has a strong derivation}, and every well-formed context is
strongly well formed.
\end{theorem}
\begin{proof}
\uses{thm:isdefeq-strong-prime, thm:ctxstrong-strong-prime, struct:orderedstrong}
\leanok
\lead{Idea} Specialise \texttt{IsDefEq.strong'}/\texttt{CtxStrong.strong'} to the three
components of \texttt{OrderedStrong}.
\lead{Caveat} The choke point through which \stSorry{} \texttt{patsStrong} reaches
\texttt{UniqueTyping.lean}, \texttt{ChurchRosser.lean}, \texttt{HeadReduction.lean} and the
Verify chapters.
\thesisref{typesys.tex thm:reg}
\end{proof}

\begin{lemma}[Level transport, ordinary]
\label{fam:isdefeq-levels}
\lean{Lean4Lean.VEnv.IsDefEq.eqUpToLevels,Lean4Lean.VEnv.IsDefEq.instL_r}
\leanok
\contrib{mixed}\srcloc{Lean4Lean/Theory/Typing/Strong.lean}{810}

\lead{In short} The right side of a definitional equality may be replaced by a level-equal
term; instantiating both sides with pointwise level-equivalent lists preserves it too.
\end{lemma}
\begin{proof}
\uses{thm:isdefeq-strong, thm:equptolevels-defeq, def:equptolevels}
\leanok
\lead{Idea} Strengthen, apply \texttt{EqUpToLevels.defeq} (resp.\ \texttt{instL} then it),
forget the annotations.
\thesisref{no thesis counterpart}
\end{proof}

\section{Typing inversions and stratified typing}

\begin{theorem}[Both sides strongly typed]
\label{thm:strong-hastype-prime}
\lean{Lean4Lean.VEnv.IsDefEqStrong.hasType'}
\leanok
\contrib{mixed}\srcloc{Lean4Lean/Theory/Typing/Strong.lean}{825}

\lead{In short} $\Gamma\Vdash e_1\equiv e_2:A$ implies $\Gamma\Vdash e_1:A$ and
$\Gamma\Vdash e_2:A$ in the syntax-directed \texttt{HasTypeStrong}.
\end{theorem}
\begin{proof}
\uses{def:hastypestrong, def:isdefeqstrong}
\leanok
\lead{Idea} Induction, rebuilding a \texttt{HasTypeStrong} derivation per side; asymmetric
rules retype the right side via \texttt{HasTypeStrong.defeq}.
\thesisref{no thesis counterpart}
\end{proof}

\begin{lemma}[Strong typing to equality]
\label{fam:hastypestrong-basic}
\lean{Lean4Lean.VEnv.HasTypeStrong.refl,Lean4Lean.VEnv.HasTypeStrong.hasType}
\leanok
\srcloc{Lean4Lean/Theory/Typing/Strong.lean}{865}

\lead{In short} A \texttt{HasTypeStrong} derivation gives the reflexive strong equality
$\Gamma\Vdash e\equiv e:A$, hence an ordinary $\Gamma\vdash e:A$.
\end{lemma}
\begin{proof}
\uses{def:hastypestrong, def:isdefeqstrong}
\leanok
\lead{Idea} Induction, mapping each structural rule to the matching \texttt{IsDefEqStrong}
congruence; \texttt{hasType} is \texttt{refl} then \texttt{defeq}.
\end{proof}

\begin{theorem}[Typing inversion lemmas]
\label{fam:strong-inversions}
\lean{Lean4Lean.VEnv.HasType.app_inv,Lean4Lean.VExpr.WF.app_inv,Lean4Lean.VEnv.HasType.lam_inv,Lean4Lean.VExpr.WF.lam_inv,Lean4Lean.VEnv.HasType.const_inv,Lean4Lean.VExpr.WF.const_inv,Lean4Lean.VEnv.HasType.bvar_inv}
\leanok
\contrib{mixed}\srcloc{Lean4Lean/Theory/Typing/Strong.lean}{886}

\lead{In short} Typing inversion, one lemma per head shape (\hl{the first three also have a
\texttt{VExpr.WF} variant}; \texttt{bvar\_inv} does not).
\begin{itemize}
\item $f\;a:V$ recovers $A,B$ with $f:\Pi A.B$ and $a:A$;
\item $\lambda A.body:V$ recovers $A$ a type and $body:B$-typed in $A::\Gamma$;
\item $\mathsf{const}\;c\;ls:V$ recovers $c$ existing with matching level arity;
\item $\mathsf{bvar}\;i:V$ recovers $i$ in scope.
\end{itemize}
\end{theorem}
\begin{proof}
\uses{thm:isdefeq-strong, thm:strong-hastype-prime, def:hastypestrong}
\leanok
\lead{Idea} Strengthen, project to \texttt{HasTypeStrong}, generalise the term and induct: only
\texttt{defeq}/\texttt{base} apply, \texttt{base} exposes the structural rule.
\lead{Caveat} All seven inherit \stSorry{} \texttt{patsStrong} (Theorem~\ref{thm:isdefeq-strong}).
\thesisref{typesys.tex thm:reg}
\end{proof}

\begin{definition}[Stratified typing]
\label{def:hastypestratified}
\lean{Lean4Lean.VEnv.HasTypeStratified,Lean4Lean.VEnv.HasTypeStratified.hasType,Lean4Lean.VEnv.HasTypeStratified.mono,Lean4Lean.VEnv.HasTypeStrong.stratify,Lean4Lean.VEnv.HasTypeStratified.to_core,Lean4Lean.VEnv.HasTypeStratified.isType}
\leanok
\uses{def:hastypestrong, def:isdefeq}
\srcloc{Lean4Lean/Theory/Typing/Strong.lean}{944}
\thesisref{unique.tex \S sec:unique}

\lead{In short} $\Gamma\vdash_n e:A$: \hl{every premise derivable at level $<n$}.
\begin{itemize}
\item sound for \texttt{HasType}, monotone in $n$;
\item every \texttt{HasTypeStrong} derivation is stratified at some level (\texttt{stratify});
\item a stratified typing has a core derivation at the same level, possibly another type;
\item the type of a term stratified at $n$ is stratified at $n-1$.
\end{itemize}
\end{definition}

\section{Simultaneous substitution}

\begin{theorem}[Substitution, strong judgment]
\label{thm:strong-substeq}
\lean{Lean4Lean.VEnv.IsDefEqStrong.substEq'}
\leanok
\contrib{mixed}\srcloc{Lean4Lean/Theory/Typing/Strong.lean}{1044}

\lead{In short} For $OrderedStrong\;env$, strong $\Gamma,\Gamma_0$ and
$W:\mathtt{Ctx.SubstEq}\,env\,U\,\Gamma_0\,\sigma\,\sigma'\,\Gamma$, a strong
$\Gamma\Vdash e_1\equiv e_2:A$ yields \hl{all three cross and same-side instances}:
\begin{itemize}
\item $\Gamma_0\Vdash e_1[\sigma]\equiv e_1[\sigma']:A[\sigma]$;
\item $\Gamma_0\Vdash e_2[\sigma]\equiv e_2[\sigma']:A[\sigma]$;
\item $\Gamma_0\Vdash e_1[\sigma]\equiv e_2[\sigma']:A[\sigma]$.
\end{itemize}
\end{theorem}
\begin{proof}
\uses{def:isdefeqstrong, fam:substeq, thm:isdefeq-strong, family:pattern-transport-iota}
\leanok
\lead{Idea} Induction, rebuilding each rule on the substituted components and composing the
left/right/cross triple; \texttt{bvar} is the one case needing \texttt{OrderedStrong}, and
\texttt{iota} rebuilds the rule via the standard transport lemmas.
\thesisref{typesys.tex thm:subst}
\end{proof}

\begin{theorem}[Substitution, definitional equality]
\label{thm:isdefeq-substdf}
\lean{Lean4Lean.VEnv.IsDefEqStrong.subst,Lean4Lean.VEnv.IsDefEq.substDF}
\leanok
\contrib{mixed}\srcloc{Lean4Lean/Theory/Typing/Strong.lean}{1273}

\lead{In short} The cross projection of \texttt{substEq'}, strong then ordinary: for
$\sigma\equiv\sigma'$ and well-formed $\Gamma,\Gamma_0$: $\Gamma\vdash e_1\equiv e_2:A$ gives
$\Gamma_0\vdash e_1[\sigma]\equiv e_2[\sigma']:A[\sigma]$.
\end{theorem}
\begin{proof}
\uses{thm:strong-substeq, thm:isdefeq-strong}
\leanok
\lead{Idea} Strengthen, apply \texttt{substEq'}, forget the annotations.
\thesisref{typesys.tex thm:subst}
\end{proof}

\begin{lemma}[Substitution corollaries]
\label{fam:subst-corollaries}
\lean{Lean4Lean.VEnv.IsDefEq.subst,Lean4Lean.VEnv.HasType.subst,Lean4Lean.VEnv.IsDefEqU.subst,Lean4Lean.VEnv.IsType.subst}
\leanok
\contrib{mixed}\srcloc{Lean4Lean/Theory/Typing/Strong.lean}{1289}

\lead{In short} One-substitution specialisations of \texttt{substDF} for typed equality,
typing, untyped equality and types.
\end{lemma}
\begin{proof}
\uses{thm:isdefeq-substdf, fam:substeq}
\leanok
\lead{Idea} Each is \texttt{substDF} with $\sigma':=\sigma$ and $W.wf$ for the source context.
\thesisref{typesys.tex thm:subst}
\end{proof}

\section{Unique typing and its consequences}

\begin{theorem}[Unique typing]
\label{thm:unique-typing}
\lean{Lean4Lean.VEnv.IsDefEq.uniq}
\leanok
\srcloc{Lean4Lean/Theory/Typing/UniqueTyping.lean}{13}

\lead{In short} \hl{A term has at most one type up to definitional equality}: for a well-formed
environment and context, $\Gamma\vdash e_1\equiv e_2:A$ and $\Gamma\vdash e_2\equiv e_3:B$ give
$\exists u,\ \Gamma\vdash A\equiv B:\mathsf{sort}\;u$. \stTainted{}
\end{theorem}
\begin{proof}
\uses{thm:isdefeq-strong, thm:strong-hastype-prime, def:hastypestratified, thm:injectivity-open}
\leanok
\lead{Idea} Strengthen, project, stratify; well-founded induction on the max of the two levels.
The application case alone appeals to \texttt{IsDefEqU.forallE\_inv\_stratified}, one of the
three open \stSorry{} lemmas of Theorem~\ref{thm:injectivity-open}.
\thesisref{unique.tex thm:unique, thm:utype}
\axfoot{sorryAx via IsDefEqU.forallE\_inv\_stratified (Injectivity.lean:21) and via VEnv.WF.patsStrong (EnvLemmas.lean:334)}
\end{proof}

\begin{theorem}[Typed and untyped agree]
\label{thm:isdefeq-iff}
\lean{Lean4Lean.VEnv.IsDefEq.uniqU,Lean4Lean.VEnv.isDefEq_iff}
\leanok
\srcloc{Lean4Lean/Theory/Typing/UniqueTyping.lean}{113}

\lead{In short} \texttt{uniqU}: the untyped form of unique typing, $\Gamma\vdash A\equiv B$.
\texttt{isDefEq\_iff}: $\Gamma\vdash e_1\equiv e_2:A$ iff both sides have type $A$ and
$\Gamma\vdash e_1\equiv e_2$ untypedly.
\end{theorem}
\begin{proof}
\uses{thm:unique-typing}
\leanok
\lead{Idea} Forwards is projection; backwards, \texttt{uniq} relates the two types and
\texttt{defeqDF} transports.
\axfoot{sorryAx via IsDefEq.uniq}
\end{proof}

\begin{lemma}[Transitivity, mixed types]
\label{fam:defeq-mixing}
\lean{Lean4Lean.VEnv.IsDefEq.trans_r,Lean4Lean.VEnv.IsDefEq.trans_l,Lean4Lean.VEnv.IsDefEq.transU_r,Lean4Lean.VEnv.IsDefEq.transU_l,Lean4Lean.VEnv.IsDefEqU.defeqDF,Lean4Lean.VEnv.IsDefEqU.of_l,Lean4Lean.VEnv.IsDefEqU.of_r,Lean4Lean.VEnv.HasType.defeqU_l,Lean4Lean.VEnv.HasType.defeqU_r,Lean4Lean.VEnv.IsType.defeqU_l,Lean4Lean.VEnv.IsDefEqU.trans}
\leanok
\srcloc{Lean4Lean/Theory/Typing/UniqueTyping.lean}{124}

\lead{In short} \hl{Composing equalities stated at different types}.
\begin{itemize}
\item compose two equalities keeping either type, or a mixed typed/untyped pair;
\item type an untyped equality from either side, or transport types along one;
\item $\vdash e_1\equiv e_2$ is transitive.
\end{itemize}
\end{lemma}
\begin{proof}
\uses{thm:unique-typing}
\leanok
\lead{Idea} Each is one line: \texttt{uniq} gives the equality of the two types,
\texttt{defeqDF} moves a derivation across it.
\thesisref{typesys.tex (1)}
\axfoot{sorryAx via IsDefEq.uniq}
\end{proof}

\begin{theorem}[Weakening is conservative]
\label{thm:isdefequ-weakn-iff}
\lean{Lean4Lean.VEnv.IsDefEqU.weakN_iff}
\leanok
\srcloc{Lean4Lean/Theory/Typing/UniqueTyping.lean}{172}

\lead{In short} For well-formed $\Gamma'$ from $\Gamma$ by inserting $n$ entries at $k$:
\hl{$\Gamma'\vdash e_1{\uparrow}\equiv e_2{\uparrow}$ iff $\Gamma\vdash e_1\equiv e_2$}.
\end{theorem}
\begin{proof}
\uses{thm:isdefeq-weakn}
\stSorry{} \lead{Missing} Left-to-right (strengthening): \texttt{sorry} at
\srcloc{Lean4Lean/Theory/Typing/UniqueTyping.lean}{174} (upstream). Right-to-left is
\texttt{IsDefEq.weakN}.
\lead{Caveat} The most widely used gap in the chapter: the whole
\texttt{weakN\_iff}/\texttt{weak'\_iff} family, hence most of \texttt{ChurchRosser.lean} and
\texttt{HeadReduction.lean}, rests on it.
\thesisref{typesys.tex thm:weak}
\end{proof}

\begin{lemma}[Conservativity of weakening]
\label{fam:weakening-iff}
\lean{Lean4Lean.VExpr.WF.weakN_iff,Lean4Lean.VEnv.IsDefEq.skips,Lean4Lean.VEnv.IsDefEq.weakN_iff',Lean4Lean.OnCtx.weakN_inv,Lean4Lean.VEnv.IsDefEq.weakN_iff,Lean4Lean.VEnv.HasType.weakN_iff,Lean4Lean.VEnv.IsType.weakN_iff,Lean4Lean.VEnv.HasType.skips,Lean4Lean.VEnv.IsDefEqU.weak'_iff,Lean4Lean.VEnv.IsDefEq.weak'_iff,Lean4Lean.VEnv.HasType.weak'_iff,Lean4Lean.VEnv.IsType.weak'_iff,Lean4Lean.VExpr.WF.weak'_iff,Lean4Lean.OnCtx.weak'_inv}
\leanok
\srcloc{Lean4Lean/Theory/Typing/UniqueTyping.lean}{177}

\lead{In short} Typed, untyped and general-renaming forms of \hl{weakening as an equivalence},
context-well-formedness inversion, and \texttt{Skips} (a judgment skipping the inserted
variables restates at a type that also skips them).
\end{lemma}
\begin{proof}
\uses{thm:isdefequ-weakn-iff, thm:isdefeq-iff}
\leanok
\lead{Idea} Each direction is \texttt{weakN} or \texttt{IsDefEqU.weakN\_iff} followed by
\texttt{uniqU}; \texttt{OnCtx.weakN\_inv} inducts on the lifting.
\thesisref{typesys.tex thm:weak}
\axfoot{sorryAx via IsDefEqU.weakN\_iff (UniqueTyping.lean:174)}
\end{proof}

\section{Church--Rosser: the abstract family of reduction rules}

\begin{definition}[Church--Rosser parameters]
\label{class:cr-params}
\lean{Lean4Lean.VEnv.Params}
\leanok
\uses{inductive:pattern-core, inductive:simple-pattern, inductive:pattern-rhs, struct:venv, def:venv-wf}
\contrib{mixed}\srcloc{Lean4Lean/Theory/Typing/ChurchRosser.lean}{12}
\thesisref{unique.tex \S sec:kappa}

\lead{In short} Fixes $env$, $env.WF$, $univs$, and an \hl{abstract family $Pat\;p\;r$ of
reduction rules} with confluence's coherence conditions.
\par\smallskip\noindent
\begin{tabular}{p{0.44\linewidth}p{0.49\linewidth}}
\toprule
Field & Meaning \\
\midrule
\texttt{pat\_simple} & redex is a \texttt{SimplePattern} \\
\texttt{pat\_uniq} & overlapping rules coincide \\
\texttt{pat\_wf} & reduct equals redex under the side conditions \\
\texttt{pat\_app\_l}, \texttt{pat\_app\_l\_uniq}, \texttt{pat\_app\_uniq} & rules do not
interfere as applications \\
\texttt{extra\_pat} & every axiom is realised by a rule \\
\texttt{pat\_env} & every registered rule is a rule (\texttt{iota}) \\
\bottomrule
\end{tabular}\par\smallskip
\lead{Caveat} The one instance, \texttt{VEnv.toParams}
(\srcloc{Lean4Lean/Theory/Typing/InductiveParams.lean}{393}), needs \texttt{DefEqsAsPats}: fails
for a \texttt{def} or the quotient rule, so this applies only to axiom-and-inductive
environments.
\end{definition}

\begin{definition}[$\iota$ rules are \texttt{Pat}]
\label{axiom:cr-params-pat-env}
\lean{Lean4Lean.VEnv.Params.pat_env}
\leanok
\uses{class:cr-params, struct:venv, rule:isdefeq-pat}
\contrib{iota}\srcloc{Lean4Lean/Theory/Typing/ChurchRosser.lean}{32}
\thesisref{no thesis counterpart}

\lead{In short} Class field $env.pats\;p\;r\to Pat\;p\;r$: \hl{the bridge letting
Church--Rosser see the rules \texttt{IsDefEq.pat} reduces by}. In the only instance,
\texttt{pat\_env := id} costs nothing since there $Pat:=env.pats$.
\lead{Caveat} The docstring at \srcloc{Lean4Lean/Theory/Typing/ChurchRosser.lean}{30} justifies
this by \texttt{extra\_pat}, but the two fields (\texttt{env.defeqs} vs.\ \texttt{env.pats}) are
independent.
\end{definition}

\begin{lemma}[Pattern plumbing]
\label{fam:cr-pattern-plumbing}
\lean{Lean4Lean.Pattern.Check.OK.weakN,Lean4Lean.Pattern.Check.OK.instN,Lean4Lean.Pattern.Matches.hasType,Lean4Lean.VEnv.HasType.matches_inv,Lean4Lean.Pattern.RHS.apply_lift',Lean4Lean.Pattern.RHS.apply_liftN,Lean4Lean.VEnv.IsDefEq.apply_pat,Lean4Lean.VEnv.NormalEq.apply_pat,Lean4Lean.VEnv.ParRed.apply_pat,Lean4Lean.VEnv.Params.pat_not_var}
\leanok
\srcloc{Lean4Lean/Theory/Typing/ChurchRosser.lean}{43}

\lead{In short} Routine transport/congruence lemmas for rule matches and right-hand sides.
\begin{itemize}
\item transport side conditions/reduct under weakening/instantiation;
\item type a matched subterm;
\item congruence of $r.apply$ in the matched subterms, per relation;
\item a bare variable is never a rule redex.
\end{itemize}
\end{lemma}
\begin{proof}
\uses{class:cr-params, inductive:pattern-matches, fam:strong-inversions, fam:defeq-mixing}
\leanok
\lead{Idea} Transport lemmas induct on the right-hand side or map over \texttt{Check.OK};
congruences induct on the \texttt{RHS} and split the spine with \texttt{HasType.app\_inv}.
\axfoot{sorryAx via HasType.app\_inv (VEnv.WF.patsStrong) and IsDefEq.uniq}
\end{proof}

\section{Normal equality}

\begin{definition}[Normal equality]
\label{def:normaleq}
\lean{Lean4Lean.VEnv.NormalEq}
\leanok
\uses{def:isdefeq, def:hastype-istype}
\srcloc{Lean4Lean/Theory/Typing/ChurchRosser.lean}{84}
\thesisref{unique.tex \S sec:kappa}

\lead{In short} $\Gamma\vdash e_1\equiv_p e_2$: \hl{everything $\kappa$ reduction cannot make
unique}. No rule for the reduction rules themselves; all $\iota$/$\delta$ content lives in
\texttt{ParRed}.
\par\smallskip\noindent
\begin{tabular}{p{0.27\linewidth}p{0.65\linewidth}}
\toprule
Rule & Meaning \\
\midrule
reflexivity & at a type \\
level equivalence & on sorts and constants \\
congruence & application (both sides typed at a common $\Pi$), lambda, $\Pi$ \\
\texttt{etaL}, \texttt{etaR} & the two directed $\eta$ rules \\
proof irrelevance & -- \\
\bottomrule
\end{tabular}\par\smallskip
\end{definition}

\begin{lemma}[Normal equality to def-eq]
\label{thm:normaleq-defeq}
\lean{Lean4Lean.VEnv.NormalEq.defeq}
\leanok
\srcloc{Lean4Lean/Theory/Typing/ChurchRosser.lean}{120}

\lead{In short} Over a well-formed context, $\Gamma\vdash e_1\equiv_p e_2$ implies the untyped
$\Gamma\vdash e_1\equiv e_2$.
\end{lemma}
\begin{proof}
\uses{def:normaleq, fam:defeq-mixing, def:isdefeq}
\leanok
\lead{Idea} Induction; lambda/$\Pi$ retype the body, $\eta$ invokes \texttt{IsDefEq.eta}, proof
irrelevance is a rule of \texttt{IsDefEq}.
\thesisref{unique.tex (3)}
\axfoot{sorryAx via IsDefEq.uniq}
\end{proof}

\begin{lemma}[Symmetry, normal equality]
\label{thm:normaleq-symm}
\lean{Lean4Lean.VEnv.NormalEq.symm}
\leanok
\srcloc{Lean4Lean/Theory/Typing/ChurchRosser.lean}{147}

\lead{In short} Over a well-formed context, $\Gamma\vdash e_1\equiv_p e_2$ implies
$\Gamma\vdash e_2\equiv_p e_1$; \texttt{etaL}/\texttt{etaR} swap.
\end{lemma}
\begin{proof}
\uses{def:normaleq, thm:normaleq-defeq}
\leanok
\lead{Idea} Direct induction, using \texttt{NormalEq.defeq} to retype in the congruence cases
and \texttt{forallE\_inv} to extend the context in the $\eta$ cases.
\thesisref{unique.tex (2)}
\axfoot{sorryAx via NormalEq.defeq}
\end{proof}

\begin{lemma}[Structural properties, normal eq]
\label{fam:normaleq-struct}
\lean{Lean4Lean.VEnv.NormalEq.weakN,Lean4Lean.VEnv.NormalEq.instN,Lean4Lean.VEnv.NormalEq.instN_r,Lean4Lean.VEnv.NormalEq.defeqDFC,Lean4Lean.VEnv.NormalEq.defeq_l,Lean4Lean.VEnv.NormalEq.weakN_inv_DFC,Lean4Lean.VEnv.NormalEq.weakN_iff}
\leanok
\srcloc{Lean4Lean/Theory/Typing/ChurchRosser.lean}{167}

\lead{In short} Normal equality is \hl{stable under weakening, substitution, context
conversion; weakening is an equivalence}.
\begin{itemize}
\item stable under weakening and substitution;
\item transport along a definitionally equal context;
\item weakening is an equivalence.
\end{itemize}
\end{lemma}
\begin{proof}
\uses{def:normaleq, thm:isdefequ-weakn-iff, fam:weakening-iff, fam:defeqdfc}
\leanok
\lead{Idea} Inductions on the derivation; \texttt{weakN\_inv\_DFC} consumes \stSorry{}
\texttt{IsDefEqU.weakN\_iff} in almost every case.
\thesisref{unique.tex item:p\_subst}
\axfoot{sorryAx via IsDefEqU.weakN\_iff}
\end{proof}

\begin{theorem}[Transitivity, normal equality]
\label{thm:normaleq-trans}
\lean{Lean4Lean.VEnv.NormalEq.trans}
\leanok
\srcloc{Lean4Lean/Theory/Typing/ChurchRosser.lean}{400}

\lead{In short} Over a well-formed context, $\equiv_p$ is transitive; with reflexivity and
symmetry, \hl{$\equiv_p$ is an equivalence relation} on well-typed terms.
\end{theorem}
\begin{proof}
\uses{def:normaleq, fam:normaleq-struct, thm:isdefeq-iff}
\leanok
\lead{Idea} Well-founded recursion on a measure, case analysis on the pairs of rules;
\texttt{etaR}/\texttt{etaL} inverts under a lifted context via \texttt{weakN\_iff}.
\thesisref{unique.tex (2)}
\axfoot{sorryAx via NormalEq.weakN\_iff and IsDefEq.uniqU}
\end{proof}

\section{Parallel reduction and confluence}

\begin{definition}[Parallel reduction]
\label{def:parred}
\lean{Lean4Lean.VEnv.ParRed}
\leanok
\uses{class:cr-params, inductive:pattern-matches, inductive:pattern-rhs}
\srcloc{Lean4Lean/Theory/Typing/ChurchRosser.lean}{473}
\thesisref{unique.tex \S sec:church\_rosser}

\lead{In short} $\Gamma\vdash e\gg e'$: \hl{Tait--Martin-L\"of parallel reduction} -- a
congruence rule per term former, $\beta$, and one \texttt{extra} rule for the abstract family.
\begin{itemize}
\item \lead{Given} $Pat\;p\;r$, $p$ matches $e$ at $m_1,m_2$, side conditions $r_2$ hold as
      \texttt{Check.OK}, every matched subterm $m_2\;a\gg m_2'\;a$;
\item \lead{Then} $e\gg r_1.apply\;m_1\;m_2'$.
\end{itemize}
\end{definition}

\begin{definition}[Complete parallel reduction]
\label{def:cparred}
\lean{Lean4Lean.VEnv.CParRed,Lean4Lean.VEnv.NonNeutral}
\leanok
\uses{def:parred, class:cr-params}
\srcloc{Lean4Lean/Theory/Typing/ChurchRosser.lean}{488}
\thesisref{unique.tex \S sec:church\_rosser}

\lead{In short} $\Gamma\vdash e\ggg e'$: as $\gg$, but \texttt{const}/application congruences
fire only when $e$ is neutral (not a $\beta$ redex, matching no rule with satisfied side
conditions). \hl{Forces every available redex, making the reduct essentially unique.}
\end{definition}

\begin{lemma}[Basic properties of $\gg$]
\label{fam:parred-struct}
\lean{Lean4Lean.VEnv.ParRed.rfl,Lean4Lean.VEnv.ParRed.weakN,Lean4Lean.VEnv.ParRed.instN,Lean4Lean.VEnv.ParRed.defeq,Lean4Lean.VEnv.ParRed.hasType,Lean4Lean.VEnv.ParRed.defeqDFC,Lean4Lean.VEnv.ParRed.weakN_inv,Lean4Lean.VEnv.CParRed.toParRed}
\leanok
\srcloc{Lean4Lean/Theory/Typing/ChurchRosser.lean}{499}

\lead{In short} Basic properties of $\gg$.
\begin{itemize}
\item reflexivity, weakening, $e_1\gg e_2 \wedge a_1\gg a_2 \Rightarrow e_1[a_1]\gg e_2[a_2]$;
\item soundness ($\Gamma\vdash e\equiv e':A$) and subject reduction;
\item transport along a definitionally equal context, inversion of weakening, $\ggg\subseteq\gg$.
\end{itemize}
\end{lemma}
\begin{proof}
\uses{def:parred, def:cparred, class:cr-params, fam:cr-pattern-plumbing, fam:strong-inversions}
\leanok
\lead{Idea} Inductions on the reduction, via \texttt{Params.pat\_wf} and
\texttt{IsDefEq.apply\_pat} for the matched-subterm congruence.
\thesisref{unique.tex thm:gg\_prop}
\axfoot{sorryAx via HasType.app\_inv and IsDefEq.uniq}
\end{proof}

\begin{theorem}[Complete reduct exists]
\label{thm:cparred-exists}
\lean{Lean4Lean.VEnv.CParRed.exists}
\leanok
\srcloc{Lean4Lean/Theory/Typing/ChurchRosser.lean}{717}

\lead{In short} If $\Gamma\vdash e:A$ over a well-formed context, \hl{a complete parallel
reduct exists}: $\Gamma\vdash e\ggg e'$ for some $e'$.
\end{theorem}
\begin{proof}
\uses{def:cparred, fam:strong-inversions}
\leanok
\lead{Idea} Structural recursion; classically decide \texttt{NonNeutral} to pick the
$\beta$/rule case or the congruence rule.
\lead{Caveat} Uses \texttt{Classical.byCases}, making confluence classical.
\thesisref{unique.tex thm:gg\_prop}
\axfoot{sorryAx via HasType.app\_inv and HasType.lam\_inv}
\end{proof}

\begin{theorem}[Triangle lemma]
\label{thm:parred-triangle}
\lean{Lean4Lean.VEnv.ParRed.triangle}
\leanok
\srcloc{Lean4Lean/Theory/Typing/ChurchRosser.lean}{766}

\lead{In short} If $\Gamma\vdash e:A$, $\Gamma\vdash e\gg e'$ and $\Gamma\vdash e\ggg o$, there
is $o'$ with $\Gamma\vdash e'\gg o'$ and $\Gamma\vdash o'\equiv_p o$. \hl{The technical core of
confluence.}
\end{theorem}
\begin{proof}
\uses{def:cparred, def:parred, def:normaleq, class:cr-params, fam:normaleq-struct}
\leanok
\lead{Idea} Nested induction, inverting the parallel reduction; the rule case uses
\texttt{pat\_uniq} to show two matching rules coincide, and $\beta$/$\eta$ closes with the
\texttt{NormalEq} congruences.
\thesisref{unique.tex thm:tri}
\axfoot{sorryAx via NormalEq and the inversion lemmas}
\end{proof}

\begin{theorem}[One-step confluence]
\label{thm:parred-church-rosser}
\lean{Lean4Lean.VEnv.ParRed.church_rosser}
\leanok
\srcloc{Lean4Lean/Theory/Typing/ChurchRosser.lean}{916}

\lead{In short} $\Gamma\vdash e:A$, $e\gg e_1$, $e\gg e_2$ give $e_1',e_2'$ with
$e_1\gg e_1'$, $e_2\gg e_2'$ and $\Gamma\vdash e_1'\equiv_p e_2'$.
\end{theorem}
\begin{proof}
\uses{thm:parred-triangle, thm:cparred-exists, thm:normaleq-trans, thm:normaleq-symm}
\leanok
\lead{Idea} Take $e$'s complete reduct, apply the triangle lemma to both sides, glue the two
normal equalities with \texttt{trans}/\texttt{symm}.
\thesisref{unique.tex thm:church\_rosser}
\axfoot{sorryAx via ParRed.triangle and NormalEq.trans}
\end{proof}

\begin{lemma}[Iterated parallel reduction]
\label{def:parreds}
\lean{Lean4Lean.VEnv.ParRedS,Lean4Lean.VEnv.ParRedS.hasType,Lean4Lean.VEnv.ParRedS.defeq,Lean4Lean.VEnv.ParRedS.defeqDFC,Lean4Lean.VEnv.ParRedS.app,Lean4Lean.VEnv.ParRedS.lam,Lean4Lean.VEnv.ParRedS.forallE,Lean4Lean.VEnv.ParRedS.inst,Lean4Lean.VEnv.ParRedS.weakN}
\leanok
\srcloc{Lean4Lean/Theory/Typing/ChurchRosser.lean}{924}

\lead{In short} $\Gamma\vdash e\gg^* e'$: \hl{the reflexive-transitive closure of $\gg$}.
Preserves typing, implies $\Gamma\vdash e\equiv e':A$, closed under the congruences,
substitution and weakening.
\end{lemma}
\begin{proof}
\uses{def:parred, fam:parred-struct}
\leanok
\lead{Idea} Induction along the closure, applying the one-step lemma at each step.
\thesisref{unique.tex \S sec:church\_rosser}
\axfoot{sorryAx via ParRed.defeq and ParRed.hasType}
\end{proof}

\begin{lemma}[Eta-expansion contexts]
\label{def:parredext}
\lean{Lean4Lean.VEnv.ParRedExt,Lean4Lean.VEnv.ParRedExt.depth,Lean4Lean.VEnv.ParRedExt.apply,Lean4Lean.VEnv.ParRedExt.meas,Lean4Lean.VEnv.IsApp,Lean4Lean.VEnv.ParRedExt.isApp,Lean4Lean.VEnv.hasType_app_bvar0}
\leanok
\srcloc{Lean4Lean/Theory/Typing/ChurchRosser.lean}{998}

\lead{In short} A context built from \texttt{lift} and ``lift then apply to
$\mathsf{bvar}\;0$'' steps -- \hl{generalises the induction over a stack of $\eta$ expansions}.
\begin{itemize}
\item \texttt{isApp}: such a context never turns an application into a non-application;
\item \texttt{hasType\_app\_bvar0}: inverts $\Gamma\vdash e{\uparrow}\;(\mathsf{bvar}\;0):V$
      into a $\Pi$ typing of $e$.
\end{itemize}
\end{lemma}
\begin{proof}
\uses{def:normaleq, thm:isdefequ-weakn-iff, fam:strong-inversions}
\leanok
\lead{Idea} \texttt{isApp} inducts on the context; \texttt{hasType\_app\_bvar0} inverts the
application, retypes by unique typing, applies $\eta$, strips the lift.
\axfoot{sorryAx via IsDefEqU.weakN\_iff}
\end{proof}

\begin{lemma}[Beta absorbed by $\equiv_p$]
\label{thm:parredext-parred-beta}
\lean{Lean4Lean.VEnv.ParRedExt.parRed_beta}
\leanok
\srcloc{Lean4Lean/Theory/Typing/ChurchRosser.lean}{1045}

\lead{In short} If $\Gamma\vdash f\equiv_p\lambda A.e'$ and $\Gamma\vdash f\,a:B$, then
$\Gamma\vdash f\,a\gg^* e$ for some $e$ with $\Gamma\vdash e\equiv_p e'[a]$.
\end{lemma}
\begin{proof}
\uses{def:normaleq, def:parredext, def:parreds}
\leanok
\lead{Idea} Structural recursion on $f$, generalised over a bounded-depth $\eta$-expansion
context; the \texttt{etaL} case recurses through the context stack.
\thesisref{unique.tex thm:gg\_compat}
\axfoot{sorryAx via NormalEq and hasType\_app\_bvar0}
\end{proof}

\begin{theorem}[Compatibility with $\gg$]
\label{thm:normaleq-parred}
\lean{Lean4Lean.VEnv.NormalEq.parRed}
\leanok
\srcloc{Lean4Lean/Theory/Typing/ChurchRosser.lean}{1181}

\lead{In short} If $\Gamma\vdash e_1\equiv_p e_2$ and $\Gamma\vdash e_2\gg e_2'$, there is
\hl{$e_1'$ with $\Gamma\vdash e_1\gg^* e_1'$ and $\Gamma\vdash e_1'\equiv_p e_2'$}.
\end{theorem}
\begin{proof}
\uses{def:normaleq, def:parred, thm:parredext-parred-beta, fam:normaleq-struct}
\stSorry{} \lead{Missing} Two upstream \texttt{sorry}s
(\srcloc{Lean4Lean/Theory/Typing/ChurchRosser.lean}{1193} and
\srcloc{Lean4Lean/Theory/Typing/ChurchRosser.lean}{1212}), where a rule step meets a
\texttt{constDF} resp.\ \texttt{appDF} normal equality -- cases the thesis settles by unique
typing and smallness of the eliminator.
\lead{Idea} Otherwise induction on the normal equality with inversion on the parallel
reduction; $\beta$ via \texttt{ParRedExt.parRed\_beta}, $\eta$/congruence via the
\texttt{NormalEq} closure lemmas.
\thesisref{unique.tex thm:gg\_compat}
\end{proof}

\begin{lemma}[Compatibility, iterated]
\label{thm:normaleq-parreds}
\lean{Lean4Lean.VEnv.NormalEq.parRedS}
\leanok
\srcloc{Lean4Lean/Theory/Typing/ChurchRosser.lean}{1286}

\lead{In short} Iterated form: $\Gamma\vdash e_1\equiv_p e_2$ and $e_2\gg^* e_2'$ give
$e_1\gg^* e_1'$ with $\Gamma\vdash e_1'\equiv_p e_2'$.
\end{lemma}
\begin{proof}
\uses{thm:normaleq-parred}
\leanok
\lead{Idea} Induction on the reduction sequence, applying \texttt{NormalEq.parRed} at each
step.
\axfoot{sorryAx via NormalEq.parRed (ChurchRosser.lean:1193, 1212)}
\end{proof}

\begin{lemma}[Church--Rosser def-eq]
\label{def:crdefeq}
\lean{Lean4Lean.VEnv.CRDefEq,Lean4Lean.VEnv.CRDefEq.normalEq,Lean4Lean.VEnv.CRDefEq.refl,Lean4Lean.VEnv.CRDefEq.defeq,Lean4Lean.VEnv.CRDefEq.symm,Lean4Lean.VEnv.CRDefEq.trans}
\leanok
\srcloc{Lean4Lean/Theory/Typing/ChurchRosser.lean}{1297}

\lead{In short} $\Gamma\vdash e_1\gg\ll e_2$: both sides typed and there are $e_1',e_2'$ with
$e_1\gg^* e_1'$, $e_2\gg^* e_2'$, $\Gamma\vdash e_1'\equiv_p e_2'$. \hl{Contains $\equiv_p$, is
an equivalence relation, and implies $\Gamma\vdash e_1\equiv e_2$.}
\end{lemma}
\begin{proof}
\uses{def:parreds, def:normaleq, thm:normaleq-parreds, fam:defeq-mixing}
\leanok
\lead{Idea} Immediate except transitivity, which runs confluence
(Theorem~\ref{thm:parreds-church-rosser}) on the two middle reducts and re-attaches the outer
equalities.
\thesisref{unique.tex thm:church\_rosser}
\axfoot{sorryAx via NormalEq.parRedS (CRDefEq.trans) and IsDefEq.uniq}
\end{proof}

\begin{theorem}[Church--Rosser theorem]
\label{thm:parreds-church-rosser}
\lean{Lean4Lean.VEnv.ParRedS.church_rosser}
\leanok
\srcloc{Lean4Lean/Theory/Typing/ChurchRosser.lean}{1302}

\lead{In short} \hl{Church--Rosser theorem}: if $\Gamma\vdash e:A$ and $e\gg^* e_1$,
$e\gg^* e_2$, then $\Gamma\vdash e_1\gg\ll e_2$ -- the two reducts converge up to normal
equality. \stTainted{}
\end{theorem}
\begin{proof}
\uses{thm:parred-church-rosser, thm:normaleq-parred, def:crdefeq}
\leanok
\lead{Idea} Double induction on the two reduction sequences, gluing one-step confluence with
compatibility at each step.
\thesisref{unique.tex thm:church\_rosser}
\axfoot{sorryAx via NormalEq.parRed}
\end{proof}

\begin{theorem}[Completeness of $\kappa$]
\label{thm:isdefeq-church-rosser}
\lean{Lean4Lean.VEnv.IsDefEq.church_rosser}
\leanok
\contrib{mixed}\srcloc{Lean4Lean/Theory/Typing/ChurchRosser.lean}{1347}

\lead{In short} $\Gamma\vdash e_1\equiv e_2:A$ implies $\Gamma\vdash e_1\gg\ll e_2$ -- \hl{the
bridge from declarative equality to reduction plus proof irrelevance}.
\end{theorem}
\begin{proof}
\uses{def:crdefeq, def:parreds, axiom:cr-params-pat-env, thm:pattern-realizes-took, class:cr-params, rule:isdefeq-pat}
\leanok
\lead{Idea} Induction: equivalence rules use the \texttt{CRDefEq} closure lemmas, congruences
the \texttt{ParRedS} congruences, $\beta$/\texttt{extra} each take one step. \texttt{iota}
converts the \texttt{Realizes} witness to \texttt{Check.OK} and takes one
\texttt{ParRed.extra} step via \texttt{Params.pat\_env}.
\thesisref{unique.tex thm:ckappa}
\axfoot{sorryAx via CRDefEq.trans and ParRedS.church\_rosser, hence NormalEq.parRed}
\end{proof}

\section{Weak head reduction, standardization and inference}

\begin{lemma}[Simple pattern shape]
\label{fam:hr-pattern-shape}
\lean{Lean4Lean.VEnv.Subpattern.varN_const,Lean4Lean.VEnv.Params.simple_app}
\leanok
\srcloc{Lean4Lean/Theory/Typing/HeadReduction.lean}{29}

\lead{In short} A subpattern of $\mathtt{varN}\;(\mathsf{const}\;c)\;n$ is again of that form;
hence any application subpattern of a rule pattern $p$ is $p$ itself.
\end{lemma}
\begin{proof}
\uses{inductive:pattern-core, class:cr-params}
\leanok
\lead{Idea} The first is induction on the subpattern relation after generalising $n$; the
second cases on \texttt{pat\_simple} and applies the first.
\end{proof}

\begin{definition}[Major premise]
\label{def:ismajorpremise}
\lean{Lean4Lean.VEnv.IsMajorPremise,Lean4Lean.VEnv.IsMajorPremise.lift',Lean4Lean.VEnv.IsMajorPremise.instN,Lean4Lean.VEnv.IsMajorPremise.lam}
\leanok
\uses{class:cr-params, inductive:pattern-matches}
\srcloc{Lean4Lean/Theory/Typing/HeadReduction.lean}{43}
\thesisref{axioms.tex}

\lead{In short} $IsMajorPremise\;e$: $e$ is matched by the function part $p_1$ of some
application subpattern $p_1\,p_2$ of a rule pattern -- \hl{the spine of a potential redex still
waiting for its major premise}. Stable under renaming and substitution; a lambda is never one.
\end{definition}

\begin{definition}[Weak head reduction]
\label{def:whred}
\lean{Lean4Lean.VEnv.WHRed,Lean4Lean.VEnv.WHRed.defeqDFC,Lean4Lean.VEnv.WHRed.weak',Lean4Lean.VEnv.WHRed.weakN,Lean4Lean.VEnv.WHRed.weakU_inv,Lean4Lean.VEnv.WHRed.parRed,Lean4Lean.VEnv.WHRed.defeq,Lean4Lean.VEnv.WHRed.hasType,Lean4Lean.VEnv.WHRed.instN}
\leanok
\uses{def:ismajorpremise, def:parred, class:cr-params}
\srcloc{Lean4Lean/Theory/Typing/HeadReduction.lean}{59}
\thesisref{axioms.tex}

\lead{In short} $\Gamma\vdash e\rightsquigarrow e'$: \hl{reduce the head of an application, the
argument once the head is a major premise, or contract a $\beta$/rule redex}
(\texttt{app}, \texttt{major}, \texttt{extra}). Contained in $\gg$, hence sound and
type-preserving; stable under weakening, renaming, substitution and context conversion.
\end{definition}

\begin{definition}[Weak head normal form]
\label{def:whnf}
\lean{Lean4Lean.VEnv.WHNF,Lean4Lean.VEnv.WHNF.bvar,Lean4Lean.VEnv.WHNF.lam,Lean4Lean.VEnv.WHNF.sort,Lean4Lean.VEnv.WHNF.forallE}
\leanok
\uses{def:whred}
\srcloc{Lean4Lean/Theory/Typing/HeadReduction.lean}{132}

\lead{In short} $WHNF\;\Gamma\;e$: no weak head reduction applies to $e$. Variables, lambdas,
sorts and $\Pi$ types always are.
\end{definition}

\begin{lemma}[Partial matches normal]
\label{thm:whnf-subpattern}
\lean{Lean4Lean.VEnv.WHNF.subpattern,Lean4Lean.VEnv.IsMajorPremise.whnf}
\leanok
\srcloc{Lean4Lean/Theory/Typing/HeadReduction.lean}{139}

\lead{In short} If a proper subpattern $p_1\ne p$ of a rule pattern matches $e$, then $e$ is in
weak head normal form; \hl{in particular, a major premise spine is}.
\end{lemma}
\begin{proof}
\uses{def:whnf, class:cr-params, fam:hr-pattern-shape, inductive:pattern-matches}
\leanok
\lead{Idea} A proper subpattern must be $\mathtt{varN}\;(\mathsf{const}\;c)\;n$; induct on the
putative head reduction, discharging cases with the \texttt{pat\_uniq} family.
\end{proof}

\begin{theorem}[Weak head determinism]
\label{thm:whred-determ}
\lean{Lean4Lean.VEnv.WHRed.determ}
\leanok
\srcloc{Lean4Lean/Theory/Typing/HeadReduction.lean}{174}

\lead{In short} $\Gamma\vdash e\rightsquigarrow e_1$ and $\Gamma\vdash e\rightsquigarrow e_2$
imply $e_1=e_2$.
\end{theorem}
\begin{proof}
\uses{def:whred, thm:whnf-subpattern, inductive:pattern-matches, class:cr-params}
\leanok
\lead{Idea} Induction on the first reduction, inversion on the second; overlapping rule redexes
excluded by \texttt{pat\_uniq} with \texttt{Pattern.matches\_determ}, head/argument overlaps by
\texttt{WHNF.subpattern}.
\end{proof}

\begin{lemma}[Iterated weak head reduction]
\label{def:whreds}
\lean{Lean4Lean.VEnv.WHRedS,Lean4Lean.VEnv.WHRedS.defeqDFC,Lean4Lean.VEnv.WHRedS.parRedS,Lean4Lean.VEnv.WHRedS.app,Lean4Lean.VEnv.WHRedS.major,Lean4Lean.VEnv.WHRedS.defeq,Lean4Lean.VEnv.WHRedS.hasType,Lean4Lean.VEnv.WHRedS.weak',Lean4Lean.VEnv.WHRedS.weakN,Lean4Lean.VEnv.WHRedS.instN,Lean4Lean.VEnv.WHNF.whRedS,Lean4Lean.VEnv.WHRedS.determ,Lean4Lean.VEnv.WHRedS.weakU_inv}
\leanok
\srcloc{Lean4Lean/Theory/Typing/HeadReduction.lean}{216}

\lead{In short} $\rightsquigarrow^*$: the closure of $\rightsquigarrow$, sound for definitional
equality, \hl{deterministic up to a normal form} (a term reduces to at most one WHNF).
\end{lemma}
\begin{proof}
\uses{thm:whred-determ, def:whnf, def:parreds}
\leanok
\lead{Idea} Inductions along the closure, applying \texttt{WHRed.determ} at each step.
\axfoot{sorryAx via WHRed.defeq and WHRed.hasType}
\end{proof}

\begin{definition}[Standard reduction]
\label{def:stred}
\lean{Lean4Lean.VEnv.StRed,Lean4Lean.VEnv.StRed.rfl,Lean4Lean.VEnv.StRed.bvar_l,Lean4Lean.VEnv.StRed.sort_l,Lean4Lean.VEnv.StRed.lam_l,Lean4Lean.VEnv.StRed.forallE_l,Lean4Lean.VEnv.StRed.defeqDFC,Lean4Lean.VEnv.StRed.parRedS,Lean4Lean.VEnv.StRed.whRed,Lean4Lean.VEnv.StRed.defeq,Lean4Lean.VEnv.StRed.weak',Lean4Lean.VEnv.StRed.weakN,Lean4Lean.VEnv.StRed.instN,Lean4Lean.VEnv.StRed.apply_pat}
\leanok
\uses{def:whreds, def:parreds}
\srcloc{Lean4Lean/Theory/Typing/HeadReduction.lean}{293}
\thesisref{no thesis counterpart (Kashima 2000)}

\lead{In short} $\Gamma\vdash e\rightsquigarrow_< e'$: \hl{weak-head-reduce first, then reduce
the immediate subterms standardly}. Contains the identity, is contained in $\gg^*$; preserved
by weakening, renaming, substitution and context conversion; four inversions read the result
off the head shape.
\end{definition}

\begin{theorem}[Standard absorbs parallel step]
\label{thm:stred-triangle}
\lean{Lean4Lean.VEnv.StRed.triangle,Lean4Lean.VEnv.StRed.triangleS}
\leanok
\srcloc{Lean4Lean/Theory/Typing/HeadReduction.lean}{409}

\lead{In short} For a context conversion $W:\Gamma_1\to\Gamma_2$: if $\Gamma_1\vdash e:A$,
$\Gamma_1\vdash e\rightsquigarrow_< e_1$ and $\Gamma_2\vdash e_1\gg e_2$, then
$\Gamma_1\vdash e\rightsquigarrow_< e_2$ (\texttt{triangleS}: the same for $\gg^*$).
\end{theorem}
\begin{proof}
\uses{def:stred, def:parred, fam:strong-inversions, fam:defeqdfc}
\leanok
\lead{Idea} Induction with inversion on the standard reduction; the rule case reassembles the
redex via \texttt{StRed.apply\_pat}.
\axfoot{sorryAx via HasType.app\_inv and HasType.lam\_inv}
\end{proof}

\begin{theorem}[Standardization]
\label{thm:parreds-standard}
\lean{Lean4Lean.VEnv.ParRedS.standard}
\leanok
\srcloc{Lean4Lean/Theory/Typing/HeadReduction.lean}{466}

\lead{In short} If $\Gamma\vdash e:A$ and $\Gamma\vdash e\gg^* e'$, then
$\Gamma\vdash e\rightsquigarrow_< e'$ -- \hl{any reduction sequence rearranges into a standard
one}. \stTainted{}
\end{theorem}
\begin{proof}
\uses{thm:stred-triangle, def:stred}
\leanok
\lead{Idea} Immediate from \texttt{StRed.triangleS} started at \texttt{StRed.rfl} with the
identity context conversion.
\thesisref{no thesis counterpart (Kashima 2000)}
\axfoot{sorryAx via StRed.triangle}
\end{proof}

\begin{theorem}[Sort head-reduction]
\label{thm:reduce-sort}
\lean{Lean4Lean.VEnv.IsDefEq.reduce_sort}
\leanok
\srcloc{Lean4Lean/Theory/Typing/HeadReduction.lean}{470}

\lead{In short} If $\Gamma\vdash e\equiv\mathsf{sort}\;u:A$, then
$\Gamma\vdash e\rightsquigarrow^*\mathsf{sort}\;u'$ for some $u'\approx u$.
\end{theorem}
\begin{proof}
\uses{thm:isdefeq-church-rosser, thm:parreds-standard, thm:isdefeq-iff, def:normaleq, thm:injectivity-open}
\leanok
\lead{Idea} Church--Rosser and standardization give a common standard reduct; case analysis on
the normal equality, using unique typing and \texttt{sort\_forallE\_inv}, forces it to be a
sort.
\thesisref{unique.tex thm:1dinv}
\axfoot{sorryAx via IsDefEq.church\_rosser and IsDefEqU.sort\_inv}
\end{proof}

\begin{theorem}[Pi head-reduction]
\label{thm:reduce-foralle}
\lean{Lean4Lean.VEnv.IsDefEq.reduce_forallE}
\leanok
\srcloc{Lean4Lean/Theory/Typing/HeadReduction.lean}{488}

\lead{In short} If $\Gamma\vdash e\equiv\Pi A.B:V$ then
$\Gamma\vdash e\rightsquigarrow^*\Pi A'.B'$ for some $A',B'$.
\end{theorem}
\begin{proof}
\uses{thm:isdefeq-church-rosser, thm:parreds-standard, thm:foralle-inv-derived, thm:injectivity-open}
\leanok
\lead{Idea} Same shape as \texttt{reduce\_sort}: Church--Rosser, standardization, then
inversion on the normal equality using unique typing and \texttt{sort\_forallE\_inv}.
\thesisref{unique.tex thm:1dinv}
\axfoot{sorryAx via IsDefEq.church\_rosser and IsDefEqU.sort\_forallE\_inv}
\end{proof}

\begin{definition}[Type inference judgment]
\label{def:infertype}
\lean{Lean4Lean.VEnv.InferType,Lean4Lean.VEnv.InferType.hasType,Lean4Lean.VEnv.InferType.determ,Lean4Lean.VEnv.InferType.weak',Lean4Lean.VEnv.InferType.weakN,Lean4Lean.VEnv.InferType.weakU_inv,Lean4Lean.VEnv.InferType.weak'_inv,Lean4Lean.VEnv.InferType.instN,Lean4Lean.VEnv.InferType.inst}
\leanok
\uses{def:whreds, def:hastype-istype}
\srcloc{Lean4Lean/Theory/Typing/HeadReduction.lean}{508}
\thesisref{axioms.tex}

\lead{In short} $\Gamma\vdash e\triangleright A$: \hl{the algorithmic reading of typing}.
\begin{itemize}
\item variables/sorts/constants read their type off the context or environment;
\item application weak-head-reduces the function's type to a $\Pi$ and checks the argument;
\item $\Pi$ reduces both components' inferred types to sorts.
\end{itemize}
Sound, deterministic, stable under renaming, weakening and substitution.
\end{definition}

\begin{theorem}[Inference completeness]
\label{thm:infertype-exists}
\lean{Lean4Lean.VEnv.InferType.exists}
\leanok
\srcloc{Lean4Lean/Theory/Typing/HeadReduction.lean}{633}

\lead{In short} If $\Gamma\vdash a:A$ then $\Gamma\vdash a\triangleright A'$ for some $A'$.
\stTainted{}
\end{theorem}
\begin{proof}
\uses{thm:isdefeq-strong, thm:strong-hastype-prime, thm:reduce-foralle, thm:reduce-sort, thm:unique-typing}
\leanok
\lead{Idea} Strengthen to \texttt{HasTypeStrong} and induct; the application/$\Pi$ cases
compare inferred and given type by unique typing, then expose the head shape with
\texttt{reduce\_forallE}/\texttt{reduce\_sort}.
\thesisref{axioms.tex; typesys.tex item:alg\_defn}
\axfoot{sorryAx twice over: via IsDefEq.church\_rosser (reduce\_sort, reduce\_forallE) and via VEnv.WF.patsStrong (IsDefEq.strong)}
\end{proof}

\begin{lemma}[Inference up to reduction]
\label{def:infertypes}
\lean{Lean4Lean.VEnv.InferTypeS,Lean4Lean.VEnv.InferTypeS.hasType,Lean4Lean.VEnv.WHRedS.inferType,Lean4Lean.VEnv.InferTypeS.determ,Lean4Lean.VEnv.InferTypeS.weak',Lean4Lean.VEnv.InferTypeS.weakU_inv}
\leanok
\srcloc{Lean4Lean/Theory/Typing/HeadReduction.lean}{659}

\lead{In short} $\Gamma\vdash e\triangleright^* A$: $e$ infers some $A'$ with
$A'\rightsquigarrow^* A$. Sound for typing, deterministic once the target is a WHNF, stable
under renaming.
\end{lemma}
\begin{proof}
\uses{def:infertype, def:whreds}
\leanok
\lead{Idea} Soundness composes \texttt{InferType.hasType} with \texttt{WHRedS.defeq};
determinism is a head induction via \texttt{InferType.determ} and \texttt{WHRedS.inferType}.
\axfoot{sorryAx via InferType.hasType}
\end{proof}

\section*{Caveats}
\label{sec:metatheory-review}

\begin{itemize}
\item \texttt{OrderedStrong} (Definition~\ref{struct:orderedstrong}) reaches \texttt{VEnv.WF}
      only through a \stSorry{} obligation; every consequence of
      Theorem~\ref{thm:isdefeq-strong} is conditional on it.
\item Two \texttt{sorry}s in \texttt{NormalEq.parRed} (Theorem~\ref{thm:normaleq-parred}) leave
      Church--Rosser, standardization and inference unproved, independently.
\item The forward direction of \texttt{IsDefEqU.weakN\_iff}
      (Theorem~\ref{thm:isdefequ-weakn-iff}) is \texttt{sorry}; most of
      \texttt{ChurchRosser.lean} and \texttt{HeadReduction.lean} depend on it.
\item Unique typing (Theorem~\ref{thm:unique-typing}) and the sort/$\Pi$-reduction lemmas
      depend on Theorem~\ref{thm:injectivity-open}, itself open.
\item The \texttt{Params} class (Definition~\ref{class:cr-params}) has one instance, for
      environments with no \texttt{def} and no quotient rule: applies to no realistic
      environment today.
\item \texttt{iota}'s \texttt{pat\_env} field (Definition~\ref{axiom:cr-params-pat-env}) is
      free there, so the \texttt{iota} case of Theorem~\ref{thm:isdefeq-church-rosser} is
      untested on a realistic environment.
\end{itemize}
</content>
